% !TEX TS-program = lualatex
Для использования программного продукта пользователь должен иметь возможность запускать его локально, что необходимо для тонкой настройки уровня сложности интеллектуальных агентов. С этой целью может быть использована система контейнеризации Docker Compose, позволяющая развернуть все компоненты приложения в изолированной и управляемой среде.

Docker Compose обеспечивает удобное управление зависимостями и конфигурациями, что особенно важно при работе с многоуровневыми системами. Локальный запуск также позволяет разработчику протестировать различные сценарии поведения агентов, оперативно вносить изменения в конфигурации и отслеживать результаты в реальном времени без необходимости публикации на удалённый сервер.

После запуска основных модулей приложения пользователь должен обучить интеллектуального агента (\refimage{edupra-build-ui}). Этот этап необходим для настройки поведения агента в соответствии с выбранными параметрами сложности и стратегии.

Для проверки и оценки сложности полученного решения можно воспользоваться модулем тестирования интеллектуальных агентов (\refimage{edupra-evaluate-ui}). Он позволяет моделировать различные сценарии взаимодействия и анализировать эффективность работы агента в различных условиях.

На \refimage{backgammon-ui-last.png} представлен итоговый интерфейс клиентского приложения. Он состоит из нескольких ключевых компонентов:
\begin{itemize}
    \item игровое поле (доска) --- отображает текущее состояние игры;
    \item терминал действий --- расположен в нижней части интерфейса и используется для ввода команд; здесь же отображается цвет активного игрока;
    \item окно действий --- находится справа и служит для упрощения ввода команд в терминал;
    \item журнал партии --- расположен слева и содержит историю всех действий, совершённых в ходе игры.
\end{itemize}

Такое разделение интерфейса позволяет пользователю эффективно управлять процессом игры, контролировать действия агентов и быстро реагировать на изменения игрового состояния.

\image[width=0.8\textwidth]{Интерфейс командной строки модуля обучения агентов}{edupra-build-ui}{edupra-build-ui}
\image[width=0.8\textwidth]{Интерфейс командной строки модуля тестирования агентов}{edupra-evaluate-ui}{edupra-evaluate-ui}
\image[width=0.8\textwidth]{Конечный вид модуля клиентского приложения}{backgammon-ui-last}{backgammon-ui-last.png}
